\chapter{Experiment: Tracing the Complete Compilation Pipeline}

\section{Objective}

This experiment will guide you through every stage of the C compilation pipeline, producing each intermediate artifact and understanding what happens at each step:
\begin{enumerate}
\item Preprocessing: source code transformation
\item Compilation: C to assembly
\item Assembly: assembly to object code
\item Linking: object files to executable
\end{enumerate}

\section{Setup Requirements}

\begin{itemize}
\item GCC or Clang compiler
\item Linux or WSL recommended
\item Basic command line familiarity
\item Tools: gcc, as, ld, nm, objdump, readelf
\end{itemize}

\section{Experiment 1: The Complete Pipeline}

\subsection{Step 1: Create a Multi-File Program}

Create a header file \texttt{calc.h}:

\begin{lstlisting}[language=C]
#ifndef CALC_H
#define CALC_H

int add(int a, int b);
int multiply(int a, int b);

#endif
\end{lstlisting}

Create \texttt{calc.c}:

\begin{lstlisting}[language=C]
#include "calc.h"

int add(int a, int b) {
    return a + b;
}

int multiply(int a, int b) {
    return a * b;
}
\end{lstlisting}

Create \texttt{main.c}:

\begin{lstlisting}[language=C]
#include <stdio.h>
#include "calc.h"

#define RESULT_VALUE 42

int main(void) {
    int x = add(10, 20);
    int y = multiply(3, 4);
    printf("add: %d, multiply: %d, result: %d\n", 
           x, y, RESULT_VALUE);
    return 0;
}
\end{lstlisting}

\subsection{Step 2: Stage 1 - Preprocessing}

Generate preprocessed output:

\begin{lstlisting}[language=bash]
gcc -E main.c -o main.i
gcc -E calc.c -o calc.i
\end{lstlisting}

Examine the preprocessed output:

\begin{lstlisting}[language=bash]
head -50 main.i
tail -30 main.i
wc -l main.c main.i
\end{lstlisting}

\textbf{What to observe:}
\begin{itemize}
\item \texttt{\#include <stdio.h>} is expanded to hundreds of lines
\item \texttt{\#include "calc.h"} is replaced with the header contents
\item \texttt{RESULT\_VALUE} is replaced with \texttt{42}
\item Header guards prevent double inclusion
\end{itemize}

\subsection{Step 3: Stage 2 - Compilation to Assembly}

Generate assembly from preprocessed source:

\begin{lstlisting}[language=bash]
gcc -S main.i -o main.s
gcc -S calc.i -o calc.s
\end{lstlisting}

Examine the assembly:

\begin{lstlisting}[language=bash]
cat main.s
cat calc.s
\end{lstlisting}

\textbf{What to observe:}
\begin{itemize}
\item Function labels (\texttt{main:}, \texttt{add:}, \texttt{multiply:})
\item Assembly directives (\texttt{.text}, \texttt{.globl}, \texttt{.type})
\item Instructions (\texttt{pushq}, \texttt{movl}, \texttt{call}, \texttt{ret})
\item References to \texttt{printf} (external symbol)
\end{itemize}

\subsection{Step 4: Stage 3 - Assembly to Object Code}

Generate object files:

\begin{lstlisting}[language=bash]
gcc -c main.s -o main.o
gcc -c calc.s -o calc.o
\end{lstlisting}

Examine object files:

\begin{lstlisting}[language=bash]
file main.o
file calc.o
nm main.o
nm calc.o
\end{lstlisting}

\textbf{What to observe:}
\begin{itemize}
\item Object files are ``relocatable''
\item \texttt{main.o} has undefined references to \texttt{add}, \texttt{multiply}, \texttt{printf}
\item \texttt{calc.o} defines \texttt{add} and \texttt{multiply}
\end{itemize}

View relocations:

\begin{lstlisting}[language=bash]
readelf -r main.o
objdump -r main.o
\end{lstlisting}

\subsection{Step 5: Stage 4 - Linking}

Link object files into executable:

\begin{lstlisting}[language=bash]
gcc main.o calc.o -o program
./program
\end{lstlisting}

Expected output:

\begin{lstlisting}
add: 30, multiply: 12, result: 42
\end{lstlisting}

\section{Experiment 2: Examining Each Artifact}

\subsection{Artifact Sizes}

Compare sizes at each stage:

\begin{lstlisting}[language=bash]
ls -l main.c main.i main.s main.o program
size main.o program
\end{lstlisting}

\textbf{Why sizes change:}
\begin{itemize}
\item \texttt{.i} file: Much larger due to header expansion
\item \texttt{.s} file: Smaller, just assembly instructions
\item \texttt{.o} file: Binary, includes symbol tables and relocations
\item Executable: Largest, includes all linked code
\end{itemize}

\subsection{Symbol Tables}

Examine symbols in object files vs executable:

\begin{lstlisting}[language=bash]
echo "=== main.o symbols ===" && nm main.o
echo "=== calc.o symbols ===" && nm calc.o
echo "=== program symbols ===" && nm program | head -20
\end{lstlisting}

\section{Experiment 3: Using Individual Tools}

\subsection{The Preprocessor (cpp)}

Run preprocessor standalone:

\begin{lstlisting}[language=bash]
cpp main.c -o main.i
gcc -E main.c -o main.i
\end{lstlisting}

\subsection{The Assembler (as)}

Assemble directly:

\begin{lstlisting}[language=bash]
as main.s -o main.o
\end{lstlisting}

\subsection{The Linker (ld)}

Link directly:

\begin{lstlisting}[language=bash]
ld main.o calc.o -o program -lc \
   --dynamic-linker /lib64/ld-linux-x86-64.so.2
\end{lstlisting}

\section{Experiment 4: Understanding the Driver}

The \texttt{gcc} command is a ``driver'' that orchestrates all tools:

\begin{lstlisting}[language=bash]
gcc -v main.c calc.c -o program 2>&1 | grep -E "cc1|as|collect2"
\end{lstlisting}

\textbf{What happens:}
\begin{enumerate}
\item \texttt{cc1} compiles C to assembly
\item \texttt{as} assembles to object files
\item \texttt{collect2} (linker wrapper) links everything
\end{enumerate}

\section{Experiment 5: Examining Dependencies}

\subsection{Include Dependencies}

View all included files:

\begin{lstlisting}[language=bash]
gcc -M main.c
gcc -MM main.c  # System headers excluded
gcc -H main.c 2>&1 | head -20
\end{lstlisting}

\subsection{Library Dependencies}

View shared library dependencies:

\begin{lstlisting}[language=bash]
ldd program
readelf -d program | grep NEEDED
\end{lstlisting}

\section{Questions for Reader}

\subsection{Basic Understanding}

\begin{enumerate}
\item List the four main stages of compilation. What does each stage produce?
\item Why is the preprocessed file (\texttt{.i}) so much larger than the original source?
\item In \texttt{main.o}, why are \texttt{add} and \texttt{multiply} listed as undefined?
\item Which tool handles macro expansion? Which handles symbol resolution?
\end{enumerate}

\subsection{Intermediate Understanding}

\begin{enumerate}
\setcounter{enumi}{4}
\item What happens when \texttt{\#include <stdio.h>} is processed?
\item What are relocation entries? Why are they needed?
\item What does the linker do that the compiler cannot?
\item Why is \texttt{gcc} called a ``driver'' rather than a compiler?
\end{enumerate}

\subsection{Advanced Understanding}

\begin{enumerate}
\setcounter{enumi}{8}
\item Why is separate compilation useful?
\item If you change only \texttt{calc.c}, which stages need to be re-run?
\item How would the pipeline differ for cross-compilation?
\item What would change in the linking stage if you used \texttt{-static}?
\end{enumerate}

\section{Summary}

Through these experiments, you should now understand:

\begin{itemize}
\item \textbf{Pipeline Stages}: Preprocessing $\rightarrow$ Compilation $\rightarrow$ Assembly $\rightarrow$ Linking
\item \textbf{Artifacts}: Each stage produces a distinct file type (\texttt{.i}, \texttt{.s}, \texttt{.o}, executable)
\item \textbf{Tool Roles}: Each tool (cpp, cc1, as, ld) has a specific responsibility
\item \textbf{Symbol Resolution}: Linker matches references to definitions
\item \textbf{Driver Orchestration}: gcc coordinates all tools automatically
\end{itemize}
